Our preliminary experiments focus on validating the system's correctness and establishing the ablation infrastructure needed for our research questions. We have not yet run full simulation-integrated experiments; the results below characterize the debate system in isolation.

\textbf{Test Suite.}
We wrote 61 unit and integration tests organized across nine test classes: model validation (8 tests), configuration (4), prompt construction and agreeableness modifiers (11), mock response generators (7), individual node functions (11), full graph execution end-to-end (6), configuration ablations (5), graph structural integrity (4), and edge cases such as missing text context and extreme agreeableness values (5). All 61 tests pass in mock mode without API access, verifying that the graph topology, state flow, conditional looping, and output parsing are correct across all configurations.

\textbf{Ablation Configurations.}
We implemented a demo runner with eight predefined configurations designed to isolate individual experimental variables (Table~\ref{tab:ablation-configs}).

\begin{table}[h]
\centering
\small
\begin{tabular}{@{}clcccl@{}}
\toprule
\textbf{\#} & \textbf{Configuration} & \textbf{Agents} & \textbf{Agree.} & \textbf{Rounds} & \textbf{Variable Tested} \\
\midrule
1 & Default baseline         & 4 & 0.3 & 1 & Baseline \\
2 & + Sentiment agent        & 5 & 0.3 & 1 & Agent count \\
3 & + Devil's Advocate       & 3+DA & 0.3 & 1 & Adversarial mode \\
4 & High agreeableness       & 3 & 0.9 & 1 & Sycophancy \\
5 & Low agreeableness        & 3 & 0.1 & 1 & Confrontation \\
6 & Three debate rounds      & 3 & 0.3 & 3 & Debate depth \\
7 & No pipeline              & 3 & 0.3 & 1 & Preprocessing \\
8 & Full system, risk-off    & 5+DA & 0.3 & 2 & Combined stress \\
\bottomrule
\end{tabular}
\vspace{4pt}
\caption{Ablation configurations for isolating experimental variables. ``Agree.'' = agreeableness knob value; ``DA'' = Devil's Advocate.}
\label{tab:ablation-configs}
\end{table}

\textbf{Preliminary Observations from Live Runs.}
We ran all eight configurations against three sample market scenarios (bullish, mixed signals, and risk-off) using \texttt{gpt-4o-mini}. Several patterns emerged that motivate our research questions:

\begin{enumerate}
\setlength\itemsep{2pt}
  \item \textbf{Causal depth varies by role.} The Macro Strategist and Value Analyst consistently produced L2 intervention claims, while the Risk Manager and Technical Analyst more often defaulted to L1 associations unless the prompt's ``MUST include L2 or L3'' constraint was active.
  \item \textbf{Agreeableness affects convergence.} Config~4 (agreeableness = 0.9) produced near-unanimous agreement in the first round with minimal revision, while Config~5 (agreeableness = 0.1) generated substantive disagreements that persisted through to the judge's final memo.
  \item \textbf{Devil's Advocate changes the judge's output.} When the adversarial agent was present (Config~3), the judge's audited memo was noticeably longer and included explicit rebuttals, compared to the non-adversarial baseline where the memo tended toward bland consensus.
  \item \textbf{Multiple rounds deepen reasoning.} Config~6 (3 rounds) showed agents upgrading their causal claims from L1 to L2 across successive revisions, consistent with the hypothesis that iterative critique improves reasoning quality.
\end{enumerate}

These observations are qualitative and based on a small number of runs. Systematic quantitative evaluation, including T\textsuperscript{3} scoring and correlation with trading metrics, is planned as part of our full experimental protocol.
